% Requires \usepackage{booktabs}
\begin{table}[!t]
  \renewcommand{\arraystretch}{1.2}
  \caption{Retune Latency by Method and Hop Type}
  \label{tab:retune-latency}
  \centering
  \small
  \begin{tabular}{lrrrr}
    \toprule
    & \multicolumn{2}{c}{In-band} & & \\
    \cmidrule(lr){2-3}
    Method & 802.11ah & 802.15.4 & Cross-band & Ratio \\
    \midrule
    Host tuning            & 10.11 (0.13) & 10.08 (0.18) & 111.63 (0.73) & 11.1$\times$ \\
    FPGA tuning            &  5.95 (0.14) &  5.91 (0.13) &  26.88 (0.16) &  4.5$\times$ \\
    Quick tune, 8 profiles &  0.30 (0.03) &  0.30 (0.04) &   0.29 (0.03) &  1.0$\times$ \\
    Quick tune, 16 profiles &  1.18 (0.04) &  1.18 (0.03) &   1.13 (0.20) &  1.0$\times$ \\
    \bottomrule
  \end{tabular}
  \par\vspace{2pt}
  \parbox{\linewidth}{\footnotesize Mean (standard deviation) in ms. Ratio is
  cross-band over in-band, i.e.\ the price of changing band. Hops to the
  already-tuned channel are excluded: they are not retunes and cost about a
  third of one. Measurements per cell: 650 / 240 / 832 for the tuning rows and
  36 / 36 / 96 and 168 / 168 / 384 for the quick-tune rows, the latter pooled
  over three independent runs.

  Host and FPGA tuning differ only in where the tuning algorithm executes; the
  FPGA path removes host round-trips and is selected with the
  \texttt{BLADERF\_DEFAULT\_TUNING\_MODE} environment variable, requiring no
  change to the application. Quick tune instead recalls a stored RF front-end
  profile, performing no tuning computation at all, which is why its cost is
  independent of hop type: changing band becomes as cheap as staying in one.
  The two quick-tune rows differ only in how many profiles are registered.
  Latency is flat up to 8 profiles and flat again from 16 onward, with 64
  profiles costing no more than 16; the mechanism behind the threshold was not
  identified. Profiles are allocated once per channel and are never recycled,
  roughly 250 allocations exhausting the device, so they must be captured at
  setup rather than per retune.}
\end{table}
